\chapter*{Abstract}
Metabolic syndrome (MetS), marked by obesity, hypertension, insulin resistance, and dyslipidemia, is a growing public health challenge, particularly among working populations. Key modifiable contributors include workplace stress, poor diet, and low physical activity. This systematic review (2015--2025), conducted using the PRISMA framework across PubMed, Scopus, and ScienceDirect, analyzed 27 eligible studies from 1,214 records. Results showed a research focus on middle-aged and older workers, with limited attention to younger populations. Physical activity was the most frequently assessed factor, while workplace stress and dietary patterns were often underexplored. Significant associations emerged between workplace type and gender, age and disease status, and gender and dietary assessments. Methodological gaps included reliance on self-reports and insufficient integration of lifestyle factors. The review concludes that multi-component workplace interventions-combining stress management, physical activity, and nutrition-are essential to reduce metabolic risk. It also calls for broader demographic inclusion, more robust assessment tools, and balanced research designs to strengthen prevention strategies.

\noindent \textbf{Keywords:} Workplace Stress, Dietary Choices, Physical Activity, Metabolic Syndrome
